\section{L'évolution des métiers et des pratiques}
La réussite technique et ergonomique de la refonte de la bibliothèque numérique ne peut se concevoir indépendamment d'une restructuration des circuits de travail humains de l'établissement. L'analyse historique et organisationnelle de la \ac{BIU Cujas} révèle que la rigidité des outils informatiques passés a non seulement figé l'interface publique, mais a également façonné, voire contraint, l'organisation des équipes internes, créant des cloisonnements.

\subsection{L'état des lieux: un double cloisonnement}
L'analyse de l'organisation interne de la \ac{BIU Cujas} révèle que les difficultés ergonomiques et éditoriales de CujasNum sont le résultat d'un double cloisonnement humain et logiciel.

Sur le plan organisationnel, la bibliothèque a été marquée par une restructuration de ses services dans le courant de l'année 2018, entraînant une importante rupture structurelle. L'un des exemples est le service informatique qui s'est scindé en deux services distincts : d'un côté, le Département d'informatique technique \ac{DIT} et de l'autre, le Département d'information documentaire \ac{DID}) \footannexe{\nameref{annexe_A_rapport_audit} pour voir l'entretien du \ac{DID}.}. Cette séparation nette entraîne un cloisonnement dans la réalisation des tâches. La spécialisation de compétence était nécessaire, mais a pour effet indirect d'éloigner les experts métiers de la gouvernance technique de la bibliothèque numérique.

Ce cloisonnement humain s'est retrouvé renforcé par l'architecture logicielle de la bibliothèque CujasNum. Le traitement documentaire et la mise en ligne fonctionnaient auparavant de façon totalement distincte : le service du patrimoine traitait des métadonnées dans le catalogue de la bibliothèque sans aucune visibilité sur le rendu final, tandis que l'implémentation des fichiers sur le serveur public s'opérait de manière automatique par des scripts en arrière-plan \footannexe{\nameref{annexe_A_rapport_audit} pour voir l'entretien avec Mebrouk Leulmi.}.

Cette opacité a transformé le back-office de CujasNum en une véritable \enquote{boîte noire} technique pour les équipes du \ac{DID} et du \ac{DPCN}. La moindre correction, dans CujasNum, demandait la sollicitation permanente du responsable du \ac{DIT} \footannexe{\nameref{annexe_A_rapport_audit} pour voir l'entretien avec Grégory Colcanap.}. Cette dépendance a ralenti le rythme de mise en valeur scientifique des fonds, mais elle a également démobilisé les agents du métier, privés de toute autonomie et de visibilité sur l'impact de leur travail dans la chaîne de production.

\subsection{Redonner l'autonomie aux métiers grâce à Codex}
Pour rompre cette dépendance technique et fluidifier les processus internes, la conception du prototype Codex introduit un changement considérable, en se basant sur un back-office collaboratif et administrable par le métier au centre de son architecture. C'était d'ailleurs l'une des attentes du \ac{DID} à propos de la refonte que \enquote{le back-office soit administrable par un bibliothécaire} \footannexe{\nameref{annexe_A_rapport_audit} pour voir l'entretien du \ac{DID}.}.

Développé en Python-Django, le prototype Codex permet la mise en place d'une gestion fine des rôles à attribuer autour des rôles métiers spécifiques et des tâches qu'ils ont à effectuer. Par exemple, \enquote{le catalogueur administre et enrichit les métadonnées descriptives en lien direct avec le \ac{SUDOC} et Alma. L'éditeur de contenu rédige les actualités, éditorialise les pages d'accueil des collections et crée des parcours thématiques} \footannexe{\nameref{annexe_A_rapport_audit} pour voir l'entretien avec Mebrouk Leulmi.}. Le technicien en contrôle qualité valide la qualité des images numérisées par les prestataires extérieurs avant leur versement et s'assure de l'avancement des différents lots.

Grâce à cette répartition des droits d'accès au sein de la partie administrative du prototype Codex, le \ac{DPCN} et le \ac{DID} auront une meilleure visibilité des différentes étapes de production à l'intérieur. Ils pourront également réaliser les micro-corrections sans dépendre constamment du responsable du \ac{DIT}, permettant une meilleure répartition du travail. Le \ac{DIT} peut s'occuper alors de la maintenance technique de la plateforme \footannexe{\nameref{annexe_A_rapport_audit} pour voir l'entretien avec Mebrouk Leulmi.}.

\subsection{La conduite du changement}
Comme le rappelle Catherine Storne-Sengel, la réussite d'un projet de réinformatisation ou de refonte documentaire ne dépend pas seulement de la performance de la nouvelle infrastructure, mais repose fondamentalement sur l'accompagnement des agents face aux mutations de leurs pratiques quotidiennes\footcite{storne-sengelAccompagnerChangementsPratique2017}.

Le passage d'un mode de production caractérisé comme artisanal à une plateforme ouverte et collaborative exige d'accompagner étroitement les équipes face à ce changement. En effet, comme le souligne le rapport du projet Orion, l'accompagnement des agents est nécessaire pour mettre en place de nouvelles pratiques de travail ou celles existantes.

Cette conduite du changement doit s'opérer autour de trois axes : former aux nouveaux standards du web sémantique, clarifier les responsabilités dans le workflow, et valorisation du contenu scientifique. En effet, les gestionnaires de données et les catalogueurs doivent être formés aux nouveaux standards du web sémantique, afin d'être capable de réutiliser des alignements sur des référentiels nationaux, comme IdRef \footcite{martinolliEnrichirSesDonnees2025}, de sorte que les notices des documents numérisés créées permettent de s'interconnecter à d'autres grâce à ce référentiel. De plus, il est important de clarifier les droits des différentes personnes dans le workflow afin que les automatisations ne les bloquent pas dans leur travail. Il faut également qu'ils puissent comprendre comment leur travail impacte le reste de la production. Le rôle du bibliothécaire évolue pour devenir un \enquote{Digital Librarian}. Comme l'a théorisé Sneerivasulu, dans les années 2000, le passage du bibliothécaire traditionnel au bibliothécaire numérique, véritable \enquote{ingénieur des connaissances chargé de créer et d'organiser des canaux de diffusion sémantiques interconnectés}.

Les bibliothécaires auront dans Codex différents espaces de valorisation pour les collections ou les actualités de la bibliothèque numérique. Ils peuvent être amené à rédiger des tutoriels afin de permettre aux chercheurs de prendre en main les nouveaux outils à leur disposition.